\chapter{Conclusions}
 \epigraph{To explain the phenomena in the world of our experience, to answer the question ``why?'' rather than only the question ``what?'', is one of the foremost objectives of all rational inquiry; and especially, scientific research in its various branches strives to go beyond a mere description of its subject matter by providing an explanation of the phenomena it investigates.}{---Carl Hempel and Paul Oppenheim, `Studies in the Logic of Explanation'}
Our standard cosmological model fits very well with the empirical data, and the values of its parameters have been able to be constrained quite accurately due to innovations in cosmological observation. The problem with this model is that it pertains to a theory---viz.\ FLRW theory, which describes the dynamical evolution of space in the RW line-element according to Einstein's relativity theory---that, for all its robustness in providing a generally applicable empirical model, when it was finally subjected to data that were precise enough to allow for meaningful discrimination in the actual values of its parameters, presented us with a picture of the Universe that has little to do with the rest of our knowledge, so that it is very different from any expectations that could be reasonably based on terrestrial theory. 

For, now, together with the problems of the Universe's unnatural appearance---i.e., that it appears to have little or no large-scale curvature, which basic FLRW theory describes as a statistical impossibility, and the fact that it appears to be isotropic on large scales, with local structures everywhere appearing similar even though they should not have been in causal contact prior to their formation,---we find ourselves fully confident that this unnatural Universe should also contain some form of dark matter, the presence of which has also been observed in local dynamics, along with some form of dark energy that is now driving the apparent acceleration of the cosmic expansion rate; and that these two exotic components currently make up $\sim96\%$ of the energy-content in the Universe. As well, we have the long-standing problem of the vacuum expectation value, which is a hundred orders of magnitude greater than the cosmological constant that we do observe. Thus, e.g., it has been hypothesised that the Universe should contain some sort of quintessence field that would have driven an inflationary epoch early on, but quickly relaxed to an asymptotic state that would resemble a cosmological constant. 

But, for all of these problems, the fact remains that the empirical model is good and well-constrained---so the issue is not one of accurately accounting for the data, but of reasonably explaining them.

Indeed, the greatest challenge for the experimentalists has not been to \textit{describe} the data, but to objectively account for the number of complexities that might be imposed upon the theory---and they've consistently found that the best model is the simplest one, with only two parameters, rather than any more complex variant.

The most common epistemological method that has been used to approach the cosmological problem is the one adapted from particle physics; viz.\ the `Shut up and calculate!' approach. However, the problem with this method is that it has no \textit{explanatory} power. In quantum theory, it \textit{has} proven to have great \textit{predictive} power, which some might even argue is better still; but even David Mermin, who likely coined the now-famous phrase, has since admonished it as being `snide and mindlessly dismissive', among other things \cite{Mermin2004}. Even so, it is without a doubt the most common practice amongst theoretical physicists today to employ this experimentalist-style method: to begin by modifying the basic theoretical apparatus or its input in some way---shut up and calculate---and see what new result comes from that, as it may be most readily interpreted within the context of a prior world-view.

The issue, though, is that the type of problem to which this method best applies is not the one that we've now got in cosmology. We've got reliable data, along with a model that beautifully accounts for them, according to a good idea of what is really going on. We've even found, parametrically speaking, that this is really about the simplest possibility that could have been: the empirical data indicate that the Universe is flat and isotropic, appearing the same even in regions that could not previously have been in causal contact (unless something like an inflationary epoch occurred), and that it has expanded, everywhere in absolute (CMB rest-frame) time, precisely as a hyperbolic sine function. The only real problem has been the lack of a good idea about how to explain these results within the context of the corresponding theory---but then one commonly takes to experimental theoretics, as a means of possibly predicting a small perturbation that would give rise to the apparent way of things.

But there is another, far more direct epistemological approach to the realistic explanation of unexpected empirical results; viz., to search for an objective re-interpretation of the theory's fundamental basis through rational speculation. The beauty of re-interpretation is that one does not even have to go to calculation, at least not to start with: the mathematical description is supposed to remain equivalent, while only the theory's meaning should change. Then, when all aspects have been rationally taken into account, one sees gross paradoxes and other significant problems simply fade away.

An example is the re-interpretation of special relativity theory that was presented in \S~\ref{sec_OJGKD}. For it may be said that Einstein did interpret the relativity of simultaneity as a noumenon: he thought that it would be most objective to say that the events that occur simultaneously relative to each observer should have \textit{really} occurred simultaneously from each individual's perspective, in their own proper frame. Then, the most straightforward interpretation that is consistent with the mathematics, is the block spacetime interpretation that he eventually believed in (\S~\ref{sec_OIS}). However, there are other options; e.g., that a noumenal present, perhaps defined as that which occurs simultaneously for each observer individually, may exist without having an objective representation; or that each observer's relative reality might be their past light cone, rather than their present simultaneous hyperplane. 

But such interpretations implicitly operate from the same interpretation of the relativity of simultaneity that was originally made by Einstein; viz.~that two events on an observer's past light cone that are described in their proper frame as having occurred simultaneously should be described by them as having \textit{really} occurred simultaneously. In contrast, I have argued for a reinterpretation of special relativity theory that would describe the relativity of simultaneity as a phenomenon within a three-dimensional universe evolving along a four-dimensional manifold with Lorentzian metrical structure, in which photons travel along the invariant null lines. In this interpretation, the noumena that occur with true simultaneity will only be described as simultaneous phenomena by observers who remain absolutely at rest. In \S~\ref{sec_OJGKD}, the situation of pens tracing out lines on a barograph sheet was used to illustrate how a particular foliation of special relativistic spacetime might be objectively described in this scenario.

With the concept of an absolute present reconciled in this objectively well-defined way (meaning that the physical description of an objective reality is equivalent in any system of cooridinates), not only is the problem of the block universe implication found to be resolved, but, since the actual events of the past and future that we commonly think of should then clearly be described as events of an \textit{ideal} spacetime formed of noumena within a real three-dimensional universe as it really evolves, the paradoxes of time travel are trivially resolved as well.

Furthermore, if the cosmic time of standard cosmology were interpreted as such an absolute time, as it may be, with a corresponding absolute rest-frame that anyhow finds empirical support in the CMB, then the same should apply. Therefore, because such paradoxes have still been thought to exist, this re-interpretation is perhaps most significant in the way that it provides for a clear distinction to be made between the real time of the universe that presently evolves, and the ideal timeline of the four-dimensional continuum of events, in the objective covariant physical description that applies as well to both.

But there remains something unsettling about the inertial and causal structures that need to be assumed in special relativity theory, which seems to be made more manifest in the proposed scenario in which the absolute time adds further structure to the spacetime manifold. A resolution to this problem was proposed in \S~\ref{sec_RPT}, where it was found that de~Sitter space alone is the maximally symmetric real Riemannian manifold which has intrinsic Lorentzian signature; thus, it was recognised that if the Einstein equation for Physical Reality did have a strictly positive cosmological constant, an appropriate causal structure would be inherent in spacetime. In particular, with an appropriate definition of absolute time to describe an objective present, it then became possible to formulate a special relativistic scenario similar to the one given in the case of flat space, but which works only from philosophically well-motivated principles.

It is therefore remarkable to note, that although relativity theory \textit{can} be formulated without an absolute time or a cosmological constant, the empirical evidence from the CMB and the success of the standard cosmological model, which suggest that there is in Reality an Absolute Time, and the evidence from the measured cosmic expansion rate, suggesting that the cosmological constant actually should be positive, really do support all of this speculation; and so, prior reasons have been found to alter our original expectations to ones which are entirely consistent with the empirical facts.

The remainder of the thesis progressed from here in roughly the same vein, to the formulation of a cosmological theory which exactly predicts the empirical observations that are made. And the key difference between this theory and the special theory of relativity that is formulated on a Minkowski background, is that although the causal structures of Minkowski space and de~Sitter space are similar, their inertial structures are in fact very different. For, when the simple definition of absolute time has been made according to the induced causal structure of a de~Sitter sphere, it follows that the paths of all particles with (non-trivial) uniform coordinate velocity in the absolute rest-frame---i.e., with non-trivial \textit{absolute motion}---are \textit{not} geodesics.

This led to the hypothesis, that the de~Sitter sphere, with its intrinsic causal and inertial structure, might actually be an \textit{absolute space} through which the universe evolves as a 3-sphere that contracts and expands in absolute time, and that inertia might be defined as a particle's tendency to remain in uniform absolute motion, so that physical mass might then be related to the fact that the worldlines of inertial particles are typically not geodesics---as opposed to the theory that gravitational mass should truly warp space so that inertial particles would follow geodesics of the resulting curved spacetime. 

This hypothesis was supposed to be supported in the following way: by recognising that in the neighbourhood of such an intertial particle it should be possible to algebraically construct a line-element that would describe a relatively symmetrical local spacetime geometry in which invariant null lines would `forever' propagate in the `radial' direction at the same coordinate velocity, and by subsequently solving for the components of the metric tensor in that coordinate system according to the requirement that it has to satisfy the vacuum Einstein equation, the local form of the statical SdS solution was recovered, and it was found that this indeed has one free parameter which should correspond directly to the mass so described. In particular, it so happens that according to this interpretation of the local SdS geometry, the mass of a particle that \textit{does} remain at rest absolutely, which travels along a fundamental geodesic of the de~Sitter sphere and may describe spacetime locally according to the statical SdS solution, is actually zero.

But then in general, it was found that there should be no objective reason to restrict the mass of a particle to positive values, but that the existence of negative mass particles should be expected to be as likely as positive masses according to the symmetry of the local SdS solution. For, since mass was supposed to be directly related to a particle's absolute motion through the cosmic 3-sphere, which is indeed parallelisable, the objective difference between positive and negative masses was interpreted as being due to both the magnitude and direction of their absolute velocity vector.

And by analysing the descriptions of the gravitational field surrounding both positive and negative masses, it was found that the two should appear to be mutually repulsive and invisible. This result was thought to be significant because it could then be interpreted as meaning: if such a universe did contain an initially homogeneous distribution of such particles, then in time these would naturally dissociate into a sort of lattice structure; and, since the expansion of such a universe would be absolute, so that it would not dynamically depend on the material content as it does in FLRW theory, the gravitational interaction between any local cluster thus formed, and an effective surrounding dark matter distribution, would be perceived locally as an additional central force. Thus, it was proposed that this repulsive dark anti-matter should be primarily concentrated in the cosmic voids that are observed in our Universe, and that the apparent void-filament large-scale structure that we observe should therefore be due to the fact that we are only able to see the luminous matter that exists.

Now, the distinction between positive, negative, and zero mass particles is also significant in that it indicates how one might algebraically construct a line-element that would be appropriate to the cosmological description from the perspective of massive particles---viz.\ as the average rest-frame of one species of particles might be defined as the bundle of null lines which all propagate in the same direction along the de~Sitter sphere, with the paths of zero-mass photons subsequently re-defined as null. The general relativistic solution to the line-element thus defined turns out to be the cosmological form of the SdS metric, which components are algebraically similar to those of the local solution, but with the radial coordinate now describing the cosmical evolution of the expanding 3-sphere.

It is important to note, that the essential difference between FLRW cosmology and the present theory may be stated finally with regards to this prescription. For, in deriving the general line-element for the background geometry in the former theory, Robertson made four basic assumptions \cite{Rugh2010}: (i.) a congruence of geodesics, (ii.) hypersurface orthogonality, (iii.) homogeneity, and (iv.) isotropy. Now, these assumptions \textit{do} apply to the proper coordinate system of comoving photon geodesics in the present theory, in which the metric does indeed have the RW form. However, although the universe is thus essentially supposed to be isotropic and homogeneous, and the diverging bundle of fundamental worldlines for massive particles do satisfy Weyl's principle, due to the assumption of an absolute time, those isotropic and homogeneous cosmic hypersurfaces are not orthogonal---in fact, they're at 45$^{\circ}$ angles---to the bundle of fundamental worldlines. Incidentally, the cosmic hypersurface is not synchronous in the proper fundamental frame, whereas this is required by Robertson's assumptions (i.) and (ii.), in accordance with the Einsteinian interpretation of the relativity of simultaneity. The other difference, that in the FLRW theory the expansion is described to occur according to the stress-energy of matter in the universe, whereas the expansion occurs as a fundamental property of the vacuum in the present theory, is really only secondary to this point, as, in either case, the restriction of the components of the metric according to Einstein's equation is subsequent to the algebraic statement of the basic line-element.  

Now, the facts, that the universe actually is isotropic and homogeneous, and that the fundamental worldlines satisfy Weyl's principle, are actually enough to meet the requirement that this universe should appear isotropic to any massive observer in the fundamental rest-frame, due to the relative co-motion of matter. But also, the factor of expansion should be constrained precisely to the rate at which the perceived cosmic radius (which was defined as the radial coordinate in the cosmological SdS solution) evolves in the proper time of fundamental observers; and, as shown in \S~\ref{sec_MC}, this is in fact identical to the flat $\Lambda$CDM scale-factor of FLRW cosmology.

Therefore, because this was all described to take place only on the expanding half of a de~Sitter sphere, beginning from a coordinate singularity at its equator, while the prior contracting half of the sphere should be appropriate to the description of the radial evolution of matter during the final stage of spherically symmetric gravitational collapse, it was proposed in \S~\ref{sec_CBHO} that our Universe, which appears to us just as the expanding universe described in this theory, might have formed through such a local gravitational collapse of matter in a prior one, with that physical process having given rise to the absolute time that needed to be assumed.







%Now, the RW metric of standard cosmology, which describes the evolution of an isotropic and homogeneous space in cosmic time, is derived through kinematical and algebraic considerations as well. But because the expansion of space in FLRW theory is then supposed to occur dynamically, according to restrictions imposed through Einstein's equation and the thermodynamical properties of the species of perfect fluid that it may contain, it is difficult to see how this might occur in the same manner as in the above interpretation of special relativity theory, since the evolving space that is described should lack the same four-dimensional background structure---the absolute space that has been supposed to really exist according to the present interpretation of special relativity theory.




